"Intro to the introduction chapter"
-- some general introduction to the chapter, maybe a brief overview of the topics covered in the chapter, and how it sets the stage for the rest of the thesis.

\subsubsection{What is Visual Programming}
Visual programming is when the programming paradigm enables the user to create programs by using visual elements rather than a text-based syntax. Instead of writing code the user will arrange visual elements like blocks, nodes, diagrams or flowcharts into programs. Doing it this way can make programming more approachables to learners that might find text-based program harder to understand\cite{MYERS199097}. Visual programming can be found in educational settings, game development and programns that have it included. 
The main idea for visual programming is to reduce complexity related to syntax. And having more focus on the logic and structure of the program, by using the visual elements of a chosen paradigm, to see and understand the relationships and execution flow of the program. 
The choice of paradigm often depends on what the user wants, as well as the type of programs they want to use when creating, and what they want to create. "graphical programming uses information in a format that is closer to the user's mental representations of problems, and will allow data to be processed in a
format closer to the way objects are manipulated in the real world" \cite{MYERS199097}

\subsection{Visual Programming Paradigms}
Visual programming languages can be categorized into differnt paradigms depending on how they represent a program's build-up and execution. Early kinds of visual programming languages describe different styles of representation, this includes flowchart and block based approaches \cite{MYERS199097}. 


Block-based visual programming represents its programs using graphical blocks that match programming constructs like loops, conditionals and expressions. The blocks are designed to join together in ways that are syntactically correct, they support this by having the blocks shape tell the user where they can be connected to another piece. This design choice supports novice programmers by having them focus on the logic rather that language rules.
Block-based systems are often used in educational settings. Scratch is a block-based programming envoirment, meant to support early programming learning with with its graphical programming elements \cite{resnick_scratch_2009}. 


Flowchart-based visual programming its programs uses diagrams to represents its structure and flow. Nodes in the diagram represent operations or decision points, and the arrow represent the flow of control between the nodes. One of the earlist visual programming paradigms was flowchart-based \cite{MYERS199097}


Node-based visual programming is often linked with dataflow systems, here the program is graphically represented as a graph that is connection to nodes. These nodes represent the program's operations, and the node's connections to other nodes represent the flow of data and execution.
Unreal Engine's blueprints uses node-based scipting language to support the interative and event-driven envoirment of Unreal Engine and its game development process \cite{unreal_blueprints}.

These paradigms are different approaches to visual programming, they share the common trait of using graphical elements to represent programming constructs, they differ in their visual representation and their usecases.

\subsection{Visual Programming in Education}
Visual programming tools have become popular in education along with the introducing computational thinking into curriculums. It has been shown that early exposure to computational competencies will benefit students in their future learning with software engineering\cite{grove_computational_2013}. Educational visual programming tools are often used to introduce the fundamentals and have reduced complexity.

\subsubsection{Why Schools Adopt Visual Programming Tools}
Schools are adopting visual programming tools to for more engaging and simpler introduction to programming, this also support the development of computational thinking skills. Computational thinking is broadly descriped as “defining patterns, generalizing from specific instances,”\cite{grover_computational_2013}. Computational thinking inlcudes a range of abilities such as pattern regognition, abstraction and algorithmic thinking. 
Tools like Scratch have been made to be an engaging and easy introduction to programming, with the ability to make games, animations or interactive stories. This allows the users to be creative and foster computational thinking skills. This can help the user learn design stratigies and problem solving skills, which can be aplied to both software engineering and nonprogramming domains\cite{resnick_scratch_2009}.




\subsection{Programming Education}
- Broader programming pedagogy
- Challenges in teaching beginners
- Teaching abstractions and architectural
- Importance of design thinking
- Importance of engagement and motivation
\subsection{Research Questions}
- Lowkey ask somebody what our research are.
- I dont know write about the workshop, that will try to answer the research questions.